By $\setN$, $\setN_+$, $\setZ$, and $\setR$ we denote the sets of nonnegative
integers, positive integers, integers, and reals, respectively.
For $i,j\in\setZ$ we put $[i,j]=\{i,i{+}1,\dots,j\}$.
The unary operation $|.|$ denotes the absolute value for numbers, the
cardinality for sets, and the length for sequences.
For a~vector $x\in\setZ^d$ ($d\in\setN$), we denote by $x(i)$ the value of its
component $i\in[1,d]$, 
hence $x=(x(1),x(2),\ldots,x(d))$. 
We use the \emph{component-wise (partial) order} $\leq$ on $\setZ^d$.


\paragraph{Vectors, restrictions.}
It will be clear from context when a vector is understood as a column vector; for example, if $B$ is an $m\times n$ matrix, then in the equation $Bx=b$ the vectors $x$ and $b$ are viewed as column vectors of dimensions $n$ and $m$, respectively. By $\mathbf{0}$ we denote the zero vector, whose dimension is clear from context.
Sometimes we consider vectors as elements of $\setZ^J$, where $J$ is a finite subset of $\setN$; for example, $\setN^d$ is identified with $\setN^{[1,d]}$.
For $x\in \setZ^d$ and $J\subseteq[1,d]$, we denote by $x\restr{J}$
the restriction of $x$ to $J$, that is, the vector in $\setZ^J$
satisfying $x\restr{J}(i)=x(i)$ for each $i\in J$. For
$X\subseteq\setZ^d$, we put $X\restr{J}=\{x\restr{J}\mid x\in X\}$.

\paragraph{Multiplication, dot product, rank.}
We use “$\cdot$” for standard multiplication (including multiplication of a vector or a matrix by a scalar), but for $x,y\in\setZ^d$ we use $\scalar{x}{y}$ to denote the dot product $\sum_{i=1}^d x(i)\cdot y(i)$.
The \emph{rank} of an $m\times n$ matrix $B$ is denoted by $\rank{B}$;
in particular, $\rank{B}\leq\min\{m,n\}$.

\paragraph{Norms for vectors, sets, and matrices.}
We use the following \emph{norms} for vectors $x\in\setZ^J$ and finite sets $X\subseteq\setZ^J$ (where $\setZ^J$ may in particular be $\setZ^d$):
\[
	\nnorm{x}=\max_{i\in J}|x(i)|, \ \ 
	\norm{1}{x}=\sum_{i\in J}|x(i)|, \ \ 
\nnorm{X}=\max_{x\in X}\nnorm{x}, \ \ 
	\norm{1}{X}=\max_{x\in X}\norm{1}{x};
\]
we stipulate that $\max\emptyset=0$.
Hence,
	\[0\leq\nnorm{x}\leq\norm{1}{x}\leq |J|\cdot\nnorm{x},
	{\textnormal{ and }}
0\leq\nnorm{X}\leq\norm{1}{X}\leq |J|\cdot\nnorm{X}.
	\]
We also define $\nnorm{B}$ and $\norm{1}{B}$ \emph{for matrices over
$\setZ$}, viewing a matrix as the set of its rows.
For an $m\times n$ matrix $B$, we have  
$\norm{}{B}=\max_{i,j}|B_{ij}|$ and
 $\norm{1}{B}=\max_{i\in[1,m]}\sum_{j=1}^n|B_{ij}|$.

\paragraph{Submononoid \bm{$X^*$}.} 
For $X\subseteq \setZ^d$, we denote by $X^*$ the \emph{submonoid of $(\setZ^d,+)$ generated by $X$}, that is, the set of all finite sums of elements of $X$.



\begin{figure}[t]
\begin{center}
\begin{tikzpicture}[node distance=1.5cm,>=stealth',bend angle=20,auto]
    \tikzstyle{place}=[circle,thick,draw=black!75,fill=blue!20,minimum size=4mm]
    \tikzstyle{transition}=[rectangle,thick,draw=black!75,
    fill=black!20,minimum size=3mm]
    \tikzstyle{dots}=[circle,thick,draw=none,minimum size=6mm]

    \tikzstyle{every label}=[black]

    \begin{scope}[
      % move tokens apart such they do not intersect with the border of the place
      token distance=5.2pt
      ]
        \node [transition] (a1) [] {$a_1$};
        \node [place,tokens=1] (p3) [label=right:$p_3$,below right of=a1] {}
        edge [post,bend right] (a1)
        edge [pre,bend left] node {2} (a1);
        \node [transition] (a4) [below right of=p3] {$a_4$}
        edge [post] (p3);
        \node [place,tokens=2] (p5) [label=right:$p_5$,above right of=p3] {}
        edge [post] (a1)
        edge [post,bend right] (a4)
        edge [pre,bend left] node {2} (a4);
        \node [place] (p4) [label=right:$p_4$,below left of=a4] {}
        edge [post] (a4);
        \node [transition] (a3) [left of=p4] {$a_3$}
        edge [pre] node[swap] {2} (p3)
        edge [post] (p4);
        \node [place] (p2) [label=left:$p_2$,left of=a3] {}
        edge [pre] (a1);
        \node [transition] (a2) [above left of=p2] {$a_2$}
        edge [pre] (p2);
        \draw[shorten >=0.3mm,shorten <=0.3mm,<->] (a2) to[in=150,out=0] ($(-1cm,-2.5cm)$) to[in=135,out=330] (p4);
        \node [place,tokens=1] (p1) [label=left:$p_1$,above right of=a2] {}
        edge [post] (a1)
        edge [pre] (a2);
        \draw[shorten >=0.3mm,shorten <=0.3mm,<->] (p1) to (a4);
    \end{scope}
    \begin{scope}[xshift=6.5cm]
        \node [transition] (a1) [] {$a_1$};
        \node [place] (p3) [label=right:$p_3$,below right of=a1] {}
        edge [post,bend right] (a1)
        edge [pre,bend left] node {2} (a1);
        \node [transition] (a4) [below right of=p3] {$a_4$}
        edge [post] (p3);
        \node [place] (p4) [label=right:$p_4$,below left of=a4] {}
        edge [post] (a4);
        \node [transition] (a3) [left of=p4] {$a_3$}
        edge [pre] node[swap] {2} (p3)
        edge [post] (p4);
        \node [place] (p2) [label=left:$p_2$,left of=a3] {}
        edge [pre] (a1);
        \node [transition] (a2) [above left of=p2] {$a_2$}
        edge [pre] (p2);
        \draw[shorten >=0.3mm,shorten <=0.3mm,<->] (a2) to[in=150,out=0] ($(-1cm,-2.5cm)$) to[in=135,out=330] (p4);
        \node [place] (p1) [label=left:$p_1$,above right of=a2] {}
        edge [post] (a1)
        edge [pre] (a2);
        \draw[shorten >=0.3mm,shorten <=0.3mm,<->] (p1) to (a4);
    \end{scope}
\end{tikzpicture}

\end{center}
\caption{A 5-dim net $A_1$ (left), and a 4-dim net $A_2$
	(right), a restriction of $A_1$.}
    \label{fig:example}
\end{figure}

\paragraph{Petri nets.}
A common definition presents a Petri net as a tuple $(P,T,W)$ where $P$ and
$T$ are finite disjoint sets of \emph{places} and \emph{transitions},
respectively, and $W\colon (P\times T)\cup(T\times P)\to \setN$
is viewed as a set of \emph{weighted arcs}.
Figure~\ref{fig:example} 
depicts places as circles, transitions as boxes, and nonzero-weight
arcs, while leaving weights equal to $1$ implicit; an arc with arrows
at both ends denotes two arcs, one in each direction. 
Nevertheless, we adopt an equivalent definition that we find notationally more convenient in our context.

\paragraph{Petri nets as sets of actions.}
A \emph{Petri net} $A$ \emph{of dimension} $d\in\setN$, a
\emph{$d$-dim net $A$} for short, is a finite set of pairs
$a=(a_-,a_+) \in \Nat^d \times \Nat^d$ 
which are referred to as 
 \emph{actions} (a synonym for \emph{transitions}).
 Each action $a=(a_-,a_+)$ has an associated \emph{displacement},
defined as $\Delta(a)=(a_+-a_-)\in\setZ^d$. 
We set 
	\[
	\nnorm{A}=\max_{a\in A}\nnorm{a},
	\textnormal{ where }
\nnorm{a}=\nnorm{\{a_+,a_-\}}.
\]
We note the close relation to \emph{vector addition
	systems} (VASs), where a Petri net action
$a=(a_-,a_+)$ is replaced by the plain vector $\Delta(a)$.

\paragraph{Markings, and relations  \bm{$\step{a}$},
\bm{$\xdashrightarrow{a}$}.}
A \emph{marking} of $A$ is a vector $x\in\Nat^d$ assigning the value $x(i)$, that is, the number of \emph{tokens}, to each \emph{component} (a synonym for \emph{place}) $i\in[1,d]$.
An \emph{action} $a$ is
\emph{enabled at a marking} $x$ if $x\geq a_-$.
The ``one-step'' \emph{relation} $\step{a}$
 is defined on $\setN^d$  by setting ${x}\step{a}{y}$
if  $x\geq a_-$ and $y=x+\Delta(a)=x-a_-+a_+$.
The \emph{virtual} ``one-step'' \emph{relation} $\virtual{a}$
 is defined on $\setZ^d$  by setting ${x}\virtual{a}{y}$
if  $y-x=\Delta(a)$.
Hence, $x \step{a} y$ implies $x \virtual{a} y$, whereas the converse, even for $x,y \in \setN^d$, does not necessarily hold.

\paragraph{Conservative nets.}
A $d$-dim \emph{net} $A$ is \emph{conservative} if there exists
$w\in(\setN_+)^d$ such that $\scalar{w}{\Delta(a)}=0$ for all $a\in A$;
in this case  ${x}\step{a}{y}$ implies $\scalar{w}{x}=\scalar{w}{y}$.
If, moreover, $w(i)=1$ for all $i\in[1,d]$, then $A$ is \emph{$1$-conservative}.
Note that in a $1$-conservative net, performing actions preserves the
total number of tokens, whereas in a conservative net, only a \emph{weighted} sum of tokens is maintained.

\begin{exa}\label{exa:basicnet}
Figure~\ref{fig:example}(left) shows
a $5$-dim net $A_1$ with $4$ actions. 
For instance,
$a_1=
\left((a_1)_-,(a_1)_+\right)=
\left((1,0,1,0,1),(0,1,2,0,0)\right)$,
$\Delta(a_1)=({-}1,1,1,0,{-}1)$,
	$\Delta(a_2)=(1,{-}1,0,0,0)$,
$\Delta(a_3)=(0,0,{-}2,1,0)$.
We have 
$(0,2,0,1,0)\step{a_2}(1,1,0,1,0)$, and
	$(0,2,0,0,0)\virtual{a_2}(1,1,0,0,0)$,
	but \emph{not}
$(0,2,0,0,0)\step{a_2}(1,1,0,0,0)$.
The net $A_1$ is conservative, as demonstrated by the vector 
$w=(1,1,1,2,1)$.
The net $A_2$ in Figure~\ref{fig:example}(right)
	is obtained from $A_1$ by removing the fifth component, that
	is, 
	place $p_5$. Note that $A_2$ is 
	not
	conservative; indeed, $\Delta(a_1a_2)=(0,0,1,0)$ precludes the
	existence of a preserved 
 weighted sum with positive component weights.
\end{exa}

\paragraph{The relation \bm{$\step{\sigma}$}, 
reachability set \bm{$\Reach(x)$}, virtual reachability relation
\bm{$\xdashrightarrow{\sigma}$}.}
Given a $d$-dim net $A$,
we define the relation
$\step{\sigma}\, \subseteq \setN^d\times\setN^d$
for any action sequence $\sigma=a_1a_2\ldots a_k$
as the
composition $\step{a_1}\circ\step{a_2}\cdots\circ\step{a_k}$.
The \emph{displacement of} $\sigma$ is given by 
$\Delta(\sigma)=\sum_{j=1}^k \Delta(a_j).$
We also refer to $x\step{\sigma}y$ as an
\emph{execution in} $A$, \emph{from} a marking $x$ \emph{to} a marking $y$.

The \emph{reachability relation in} $A$ 
is $\reach\,\subseteq\setN^d\times\setN^d$, where
$x\xrightarrow{*}y$ holds if $x\xrightarrow{\sigma}y$ for some
action sequence $\sigma$. The \emph{reachability set of} a marking $x$
is defined as $\Reach(x)=\{y\mid x\reach y\}$.
We also define the relation $\virtual{\sigma}\, \subseteq
\setZ^d\times\setZ^d$ by setting $x\virtual{\sigma}y$ if
$y-x=\Delta(\sigma)$,
and we refer to  $x\virtual{\sigma}y$ as a \emph{virtual execution in $A$}.
The \emph{virtual reachability relation in} $A$ 
is $\vreach\,\subseteq\setZ^d\times\setZ^d$, where
$x\vreach y$ holds if $x\virtual{\sigma}y$ for some
action sequence $\sigma$.
 We set
\[A\D=\{\Delta(a)\mid a\in A\},\]
and note that 
 $x\virtual{*} y$ if and only if $(y-x)$ belongs to the submonoid $(A\D)^*$ of
 $(\setZ^d,+)$.
We further write $A\D^*$ instead of $(A\D)^*$.

\paragraph{Reversible nets.}
A~net $A$ is 
\emph{structurally reversible}, referred to as \emph{reversible} in this
paper, if the monoid $A\D^*$ is a~subgroup of $(\setZ^d,+)$, that is, if for every
$a\in A$ we have $-\Delta(a)\in A\D^*$.
Note that $A$ is reversible if and only if
 virtual reachability is symmetric ($x\virtual{*}y$ iff
$y\virtual{*}x$).

\paragraph{Restricted nets and their executions.}
For a $d$-dim net $A$ and $I\subseteq [1,d]$, 
\begin{equation*}
	\textnormal{the $|I|$-dim net $A\restr{I}$}
\end{equation*}
	is the \emph{restriction of} $A$ \emph{to} the
	\emph{components} (with indices)  in $I$.
	Each action $a=(a_-,a_+)$ of $A$
	gives rise to an action
	$a\restr{I}=({a_-}\restr{I},{a_+}\restr{I})$ of $A\restr{I}$.
We refer to 
executions 
$x\step{\sigma\restr{I}}y$ in $A\restr{I}$,  implicitly assuming 
that $x,y\in\setN^I$, and
adopt the convention that
\begin{equation*}
	\textnormal{$x\restr{I}\step{\sigma}y\restr{I}$ stands for
	$x\restr{I}\step{\sigma\restr{I}}y\restr{I}$.} 
\end{equation*}
Analogous conventions apply to virtual executions
$x\virtual{\sigma\restr{I}}y$, 
$x\restr{I}\virtual{\sigma}y\restr{I}$.

\begin{rem}\label{rem:dimzero}
Our definition does not exclude the nets of dimension $0$. Though these
trivial nets are of no interest, they may arise in our context as
$A\restr{I}$ when $I=\emptyset$. Formally, the empty tuple
	$()$ is the unique marking of the net $A\restr{\emptyset}$,
	and we have $()\step{a}()$ for all actions $a\in A$.
\end{rem}	

\begin{prop}[\textbf{Restrictions preserve reversibility}]\label{prop:restrpresrev}
\hfill\\
If a $d$-dim net $A$ is reversible and $I\subseteq[1,d]$, then 
$A\restr{I}$ is reversible.
\end{prop}

% Fix Indentation
\noindent
Note that this implication does not hold for
conservativeness, as demonstrated by Example~\ref{exa:basicnet}, where
$A_1$ is conservative but its restriction $A_2=(A_1)\restr{[1,4]}$ is not.


\paragraph{Liveness, structural liveness, dead markings.}
A \emph{marking} $x\in\setN^d$ of a $d$-dim net $A$ is \emph{live} if
for every action $a\in A$ and every $x'\in \Reach(x)$, 
 there exists $y\in \Reach(x')$ such that
 $y\geq a_-$ (that is, $a$ is enabled at $y$, allowing the step
 $y\step{a}y+\Delta(a)$).  A~ $d$-dim \emph{net} $A$ is \emph{structurally
 live} if there exists a marking $x\in\setN^d$ of $A$ that is live.
 Note the following standard fact:

\begin{prop}[\textbf{Liveness is closed under
	reachability}]\label{prop:liveunderreach}
\hfill\\	
If a marking $x$ of a net $A$ is live, then all
 markings $y\in\Reach(x)$ are also live.
\end{prop}

% Fix Indentation
\noindent
An \emph{action} $a\in A$ is \emph{dead at} $x$ if 
no marking in $\Reach(x)$ enables
$a$; that is, for each $y\in
\Reach(x)$ we have $y\ngeq a_-$.
For technical convenience, we say that
a \emph{marking} $x$ 
is \emph{dead} if some \emph{action} $a\in A$ is \emph{dead at} $x$.
(While in other contexts a marking is typically considered dead only if
\emph{all} actions are
dead, we adopt this weaker notion throughout the paper.) 
With this terminology, we observe the following simple fact:

\begin{prop}[\textbf{Nonliveness as reachability of dead markings}]\label{prop:nonlivereachdead}
\hfill\\
	A marking $x$ of a net $A$ is nonlive if and only if 
	$\Reach(x)$ contains a dead marking.
\end{prop}



\begin{exa}\label{exa:basiclive}
The net $A_1$ in Figure~\ref{fig:example} 
is structurally live. Indeed, $x_0=(1,0,1,0,2)$ is a
    live marking, which is easily verifiable since there is only one (infinite)
    execution from $x_0$: $(1,0,1,0,2) \step{a_1} (0,1,2,0,1) \step{a_3}
    (0,1,0,1,1) \step{a_2} (1,0,0,1,1) \step{a_4}
    (1,0,1,0,2)\step{a_1}\cdots$.

   By adding a token in $p_4$, we obtain a nonlive marking
	$x'_0=(1,0,1,1,2)$. Specifically,
	$x'_0\step{a_1a_2a_1}y=(0,1,3,1,0)$,
where $y$ is dead because actions $a_1$ and $a_4$ are dead at $y$
	(due to
	$y(5)=0$).
We could also verify that a necessary condition for $x$ to be live is
that 
$x(i)< 3$ for all $i\in[1,4]$; however, there is no such bound on $x(5)$.
(If $x(i)\geq 3$ for some $i\in[1,4]$, then $x$ can reach a~dead
	marking; on the other hand, there exist live markings $x$ with arbitrarily
	large $x(5)$.)
\end{exa}

\begin{rem}
Our example also demonstrates that liveness is not monotonic, even
in conservative nets; that is, $x\leq x'$ and $x$ being live does not
	necessarily imply that $x'$
is live. In general, the set of live markings of
	a $1$-conservative net is not even semilinear (that is, not expressible
	in Presburger arithmetic); the example
	from~\cite{DBLP:journals/acta/JancarP19} can be adapted to
	$1$-conservative nets.
\end{rem}	


\paragraph{Linear Systems \bm{$S$}, and their size-parameters
\bm{${\nnorm{S}}$},
\bm{${\mlcm{S}}$}.} 
We now define systems that are, in essence, 
quantifier-free formulas of Presburger arithmetic with divisibility 
constraints. Such a formula is $d$-dimensional if it contains $d$ 
integer variables.

Let $\bvar{x}$ be a variable vector ranging over $\setZ^d$, a \emph{$d$-dim
variable} for short, viewed as a tuple
$\bvar{x}=(\bvar{x}_1,\bvar{x}_2,\dots,\bvar{x}_d)$ of $d$ integer
variables.
A constraint of
the form $\scalar{\alpha}{\bvar{x}}\sim c$ where $\alpha\in \setZ^d$,
$c\in\setZ$, and $\sim\mathop{\in}\,\{=,\geq\}$ is an \emph{equality
constraint} if $\sim$ is $=$, and an \emph{inequality constraint} if
$\sim$ is $\geq$. It is a~\emph{homogeneous constraint} if $c=0$.
For $m\in\setN_+$, a~constraint 
$\scalar{\alpha}{\bvar{x}}\equiv c\,(\bmod\, m)$, where
$\alpha\in[0,m{-}1]^d$ and $c\in [0,m{-}1]$, is a \emph{divisibility constraint}.

A \emph{$d$-dim linear system} $S$
is a propositional formula whose atomic propositions are 
equality, inequality, and divisibility constraints for a fixed
$d$-dim variable.
The \emph{set of solutions} $\sem{S}$ 
 consists of the vectors $x\in\setZ^d$
satisfying $S$; if $\sem{S}\neq\emptyset$, then 
$S$ is \emph{satisfiable}. 

By the \emph{norm}
$\nnorm{S}$ we mean the least $s\in\setN$ such that 
$\max\{\nnorm{\alpha},|c|\}\leq s$
for all equality and inequality constraints $\scalar{\alpha}{\bvar{x}}\sim c$ in
$S$.
Moreover, $\mlcm{S}$ is the \emph{least common multiple} of all $m$
occurring in $(\bmod\, m)$ in 
divisibility constraints 
in $S$, with $\mlcm{S}=1$ if no such
constraints exist.

\paragraph{Small solution theorems.}
When dealing with solutions of linear systems, \emph{small-solution
theorems} provide bounds on the size of minimal solutions. 
Given a set $M\subseteq \setN^d$, 
by $\min_\leq(M)$ we denote the set of minimal vectors in $M$ with
respect to
the component-wise order $\leq$. Since $\leq$ is 
a~well-partial-order on $\setN^d$, the set $\min_\leq(M)$ is
always finite.
In~\cite{DBLP:conf/rta/Pottier91}, 
Pottier provided several
bounds for minimal
nonnegative solutions of a conjunction of homogeneous equality constraints. 
We now recall a bound that is particularly useful for our purposes.


\begin{lemC}[\cite{DBLP:conf/rta/Pottier91} (\textbf{Homogeneous
	equality constraints})]\label{lem:smallsolution}
\hfill\\	
  Let $M=\{x\in\setN^n \mid Bx=\mathbf{0}\}$ where $B$ is an $m\times n$
	matrix over $\setZ$. Then
	$X=\min_{\leq}(M\smallsetminus\{\mathbf{0}\})$ is a finite set
	such that $M=X^*$. Moreover the following bound holds, where
	$r=\rank{B}$:
\begin{equation*}
\norm{1}{X}\leq (1+\norm{1}{B})^r.
\end{equation*}
\end{lemC}

% Fix Indentation
\noindent
Note that the rank bound $r$ in
Lemma~\ref{lem:smallsolution} satisfies $r\leq\min\{m,n\}$.
There is also a~bound on the
solutions of a conjunction of inequality constraints, 
$Bx\geq b$, in~\cite{DBLP:conf/rta/Pottier91}, but it involves an exponent
of $m$. This is not convenient for us when $m$, the number of
constraints, is much larger than $n=d$, the dimension of our problem.
Therefore in Section~\ref{sec:smallsolutions}, we provide a proof of
the following theorem:


\begin{thm}[\textbf{Bounds on solutions of linear systems}]\label{thm:smallsolutions}
\hfill\\
	Every satisfiable $d$-dim linear system $S$ 
 has a solution $x\in\setZ^d$ such that
	\[\nnorm{x}\leq
		\mlcm{S} \cdot \left(1+ (d+1)\,!\cdot (d\cdot
		\nnorm{S}+1)^{d}\right).\]
\end{thm}

\begin{rem}
Later, in
the proof of Theorem~\ref{th:upper}, we employ a $d'$-dim linear system $S$
with $d'=d\cdot2^d$. 
It is crucial for our purposes that Theorem~\ref{thm:smallsolutions} then yields a solution that is
	at most doubly exponential (2-exp) in $d$.  
In contrast, using~\cite[Theorem 3.12]{DBLP:journals/tocl/Klaedtke08},
	we could only derive
	that for $S$ (even without divisibility constraints),  
the size of the minimal automaton encoding $\sem{S}$ 
in binary is bounded by $(2+2\cdot\nnorm{S})^{|S|}$,
where $|S|$ is the number of constraints in $S$. In our case, $|S|\geq
2^d$, and thus the bounds on the minimal solutions of $S$ derived from the
shortest accepting paths in such an automaton would be at least
$2^{(2+2\cdot\nnorm{S})^{|S|}}$,
	which is triple-exponential (3-exp) in $d$.
\end{rem}



\paragraph{Linear Systems for Subgroups of \bm{$\setZ^d$}.}
In the sequel, by a~\emph{group} we mean 
a \emph{subgroup} of $(\setZ^d,+)$,
which is also called a \emph{lattice} in this context.
The group \emph{spanned} by a set $X\subseteq \setZ^d$ is 
the monoid $(X\cup-X)^*$, that is, the set of finite integer linear
combinations of elements of $X$. 
In Section~\ref{sec:virtual}, we prove the following theorem, which
provides a way to encode any such group by a linear system, with a
bound on its  size.




\begin{thm}[\textbf{Groups expressed by linear systems of bounded size}]\label{thm:virtual2system}
\hfill\\	
Let $L$ be the group spanned by a finite set $X\subseteq
\setZ^d$. There exists a~$d$-dim linear system $S$ such that
\begin{center}
$\sem{S}=L$ and $\max\{\nnorm{S},\mlcm{S}\}\leq d\,!\cdot
	\norm{}{X}^d$.
\end{center}
\end{thm}

% Fix Indentation
\noindent
We thus obtain a corollary characterizing the virtual reachability of
reversible nets, recalling that $A\D=\{\Delta(a)\mid a\in A\}$: 
\begin{cor}[\textbf{Reversible virtual reachability expressed by a linear system}]\label{cor:virtual2system}
\hfill\\
	For every reversible $d$-dim net $A$, there exists
	a~$2d$-dim linear system $S_A$ such that
		\begin{center}
		$\sem{S_A}=\{(x,y)\in\setZ^d\times\setZ^d \mid
			x\virtual{*} y\}$, and
	$\max\{\nnorm{S_A},\mlcm{S_A}\}\leq d\,!\cdot\nnorm{A\D}^d$.
\end{center}
\end{cor}
\begin{proof}
	We recall that $x\virtual{*} y$ iff $y-x\in A\D^*$, and that 
	$A\D^*$ is a subgroup of $(\setZ^d,+)$ since $A$ is
reversible.
	Let $S$ be the linear system guaranteed 
	by 
	Theorem~\ref{thm:virtual2system} for the group $L=A\D^*$
	(spanned by the set $A\D$)
	 with a~$d$-dim variable
	 $\bvar{x}'=(\bvar{x}'_1,\bvar{x}'_2,\dots,\bvar{x}'_d)$.

	We construct $S_A$ with a~$2d$-dim variable 
	$(\bvar{x}_1,\bvar{x}_2,\dots,\bvar{x}_d,\bvar{y}_1,\bvar{y}_2,\dots,\bvar{y}_d)$ 
	 by replacing every occurrence of
	$\bvar{x}'_i$ in $S$ with the expression
	$(\bvar{y}_i-\bvar{x}_i)$. It follows that 
	\[\sem{S_A}=\{(x,y)\in\setZ^d\times\setZ^d\mid
	y-x\in\sem{S}\},\]
	which is precisely the set of pairs $(x,y)$ such that $x\virtual{*} y$.
Since the transformation only involves substituting a difference of
	variables for a single variable, the norms and the lcm remain
	unchanged: $\nnorm{S_A}=\nnorm{S}$, $\mlcm{S_A}=\mlcm{S}$.
 The claimed bound $d\,!\cdot\nnorm{A\D}^d$
then follows directly from
	Theorem~\ref{thm:virtual2system}.
\end{proof}
% Fix Indentation
\noindent
We will employ Theorem~\ref{thm:smallsolutions} and
Corollary~\ref{cor:virtual2system} to prove our main result,
Theorem~\ref{th:upper}. Before stating the result, we introduce
some further relevant notions.

\paragraph{Live bottom SCCs in reachability graphs.}
Given a $d$-dim net $A$, we denote by $(\setN^d,\to)$ 
the \emph{reachability graph} of
$A$, where $x\to y$ if $x\step{a}y$
for some $a\in A$. A nonempty set $X\subseteq \setN^d$ such that
$\Reach(x)=X$ for all $x\in X$ corresponds to 
 a \emph{bottom strongly connected component}  
in the graph; we simply call $X$ a \emph{bottom SCC} of $A$ in this case. 
A bottom SCC $X$ is \emph{live} if some $x\in X$ is live (which
implies that all $x\in X$ are live, by
Proposition~\ref{prop:liveunderreach}).

\paragraph{1-exp and 2-exp functions.}
We call a function $f\colon\setN^2\to\setN$ 
a \emph{1-exp} function if 
\[f(m,i)\leq (2+m)^{\text{poly}(i)},\]
meaning that there exists a polynomial $p\colon\setN\to\setN$
such that 
$f(m,i)\leq (2+m)^{p(i)}$ for all $m,i\in\setN$.
It is a \emph{2-exp} function if 
\[f(m,i)\leq (2+m)^{2^{\text{poly}(i)}}.\]
Throughout the paper, we specifically use values of the form $f(\nnorm{A},d)$ or 
$f(\nnorm{A},|I|)$ for $I\subseteq[1,d]$, where $A$ is a $d$-dim net.

\begin{thm}[\textbf{Bounds on minimal live markings in 
	nets with live bottom SCCs}]\label{th:upper}
	There exists a 2-exp function $f$ such that 
	for every $d$-dim net $A$ possessing a live bottom SCC, 
	there exists a live marking
	$x\in\setN^d$
	satisfying $\nnorm{x}\leq f(\nnorm{A},d)$.
\end{thm}

\paragraph{Live bottom SCCs, reversibility, conservativeness,
	structural boundedness.}
 We say that a \emph{marking} $x\in\setN^d$ of a~$d$-dim net $A$ is
 \emph{bounded} if the set of reachable markings $\Reach(x)=\{y\mid
 x\reach y\}$ is finite; that is, 
 there exists $B_x\in\setN$ such that $\nnorm{y}\leq B_x$ for
 all $y\in\Reach(x)$. A \emph{net} $A$ is \emph{structurally
 bounded} if every marking $x$ of $A$ is bounded.

The following lemma shows that structurally live conservative nets coincide
	with structurally live structurally bounded nets, and that 
	Theorem~\ref{th:upper} applies to these nets.

\begin{lem}[\textbf{Conservativeness vs.~structural boundedness for
	live nets}]\label{lem:consstrbound}\hfill
	\begin{enumerate}
\item
Every conservative net is structurally bounded.
\item
Every structurally bounded net that is structurally live has a live
			bottom SCC.	
		\item
Every net possessing a live bottom SCC is reversible.			
		\item
Every structurally bounded net that is reversible is conservative.
\end{enumerate}
\end{lem}
\begin{proof}
	(1)
	If a $d$-dim net $A$ is conservative, there exists $w\in(\setN_+)^d$
	such that $\scalar{w}{\Delta(a)}=0$ for all $a\in A$. It
	follows that 
	$x\reach y$ implies $\scalar{w}{x}=\scalar{w}{y}$. Since  
 all components of $w$ are positive, the set $C_x=\{y\in\setN^d\mid
	\scalar{w}{y}=\scalar{w}{x}\}$ is finite for any $x\in\setN^d$. Since
	$\Reach(x)\subseteq C_x$, the net $A$ is structurally bounded.

(2)
Let $x$ be a live marking of a structurally bounded net $A$.
	Since $x$ is bounded, $\Reach(x)$ is finite, which implies
	that there must
	exist a (finite) bottom SCC $X\subseteq\Reach(x)$ in the
	reachability graph of $A$.
	Since $x$ is live, every $y\in\Reach(x)$ is also live (by
	Proposition~\ref{prop:liveunderreach});
	consequently, $X$ is a live bottom SCC of $A$.

	(3) Let $X$ be a (finite or infinite) live bottom SCC of a net
	$A$.  By the definition of 
liveness,
	for
 each action $a\in A$, there exists some $y\in X$ at which $a$ is enabled,
that is, 
 $y\step{a}y+\Delta(a)$. Since $X$ is a bottom SCC, we must have
	$y+\Delta(a)\reach y$.
	This implies that $y-(y+\Delta(a))=-\Delta(a)\in A\D^*$. 
	Therefore, $-\Delta(a)\in A\D^*$ for every $a\in A$, which
	shows that $A$ is
	reversible.

	(4)
	Let $A=\{a_1,a_2,\ldots,a_k\}$ be a $d$-dim net 
that is structurally bounded and reversible.
	For each $i\in[1,d]$, let $e_i\in\setN^d$
	be the $i$th unit vector (satisfying
	$e_i(i)=1$ and $e_i(j)=0$ for all $j\neq i$). We observe that
	there exists no
	$v\in \setN^k$ such that
	\begin{equation}\label{eq:biggerei}
		\sum_{j=1}^{k} v(j)\cdot\Delta(a_j) \;\geq\; e_i.
	\end{equation}
Indeed, if such $v$ existed, then for an action sequence 
$\sigma$ with Parikh vector $v$  
	(\eg, $\sigma=(a_1)^{v(1)}\ldots (a_k)^{v(k)}$)
	and a sufficiently large marking $x$, we would have
	 $x\step{\sigma}x+\Delta(\sigma)$ with $\Delta(\sigma)\geq
	 e_i$.
This would imply $x\step{\sigma^n}x+n\cdot\Delta(\sigma)\geq x+n\cdot
	e_i$ for all $n\in\setN$, contradicting structural
	boundedness.

Since $A$ is reversible, every $\Delta(a_j)$ can be ``undone'' by
	some action sequence; that is, $-\Delta(a_j)=\sum_{\ell=1}^k
	u(\ell)\cdot\Delta(a_\ell)$ for some $u\in\setN^k$.
Hence every integer combination of the vectors $\Delta(a_j)$ equals some
	nonnegative one, so \eqref{eq:biggerei} has no solution
	$v\in\setZ^k$ either.
	Finally, the $\Delta(a_j)$ are integer vectors, so if
	\eqref{eq:biggerei} has a real solution then it has a rational one,
	and hence an integer one after multiplying by a common denominator.
	Therefore \eqref{eq:biggerei} has no solution $v\in\setR^k$ either.

Now we recall the following variant of Farkas' lemma 
(see, \eg, Proposition 6.4.3(iii)
in~\cite{DBLP:books/daglib/0016926}):
\begin{quote}
	Let $M$ be a real matrix with $m$ rows and $n$ columns, and let
	$b \in \setR^m$ be a vector.
	The system $Mx \leq b$ has a solution $x \in \setR^n$ if and only if
	every nonnegative $y \in \setR^m$ with $y^{\transp} M =
	\mathbf{0}^{\transp}$ also satisfies $y^{\transp} b \geq 0$.
\end{quote}
Fix $i\in[1,d]$, let $M$ be the $d\times k$ matrix whose $j$th column
is $-\Delta(a_j)$, and let $b=-e_i$.  Then $Mv\leq b$ is
precisely~\eqref{eq:biggerei}, which has no solution $v\in\setR^k$.
By the lemma, there exists a nonnegative
$y\in\setR^d$ with $y^{\transp}M=\mathbf{0}^{\transp}$ and
$y^{\transp}b<0$.  The $j$th entry of $y^{\transp}M$ is
$\scalar{y}{-\Delta(a_j)}$, so the former condition amounts to
$\scalar{y}{\Delta(a_j)}=0$ for all $j\in[1,k]$; and
$y^{\transp}b=\scalar{y}{-e_i}=-y(i)$, so the latter amounts to
$y(i)>0$.
By the rationality argument above, $y$ may be chosen to be a rational
vector, and hence an integer vector after multiplying by a common
denominator; both conditions are preserved.
Writing $w_i$ for the resulting vector, we
obtain for each $i\in[1,d]$ some $w_i\in\setN^d$ with $w_i(i)>0$ and
$\scalar{w_i}{\Delta(a_j)}=0$ for all $j\in[1,k]$.

	Setting $w=\sum_{i=1}^d w_i$ yields a vector $w\in(\setN_+)^d$
	satisfying $\scalar{w}{\Delta(a_j)}=0$ for all $j\in[1,k]$.
	Consequently, the net $A$ is conservative.
\end{proof}


% Fix Indentation
\noindent
The EXPSPACE-hardness result  
from~\cite{DBLP:journals/acta/JancarP19} will be slightly
strengthened by Theorem~\ref{thm:lowerbound} in
Section~\ref{sec:hardness}; Theorem~\ref{th:upper} thus allows us to
obtain a matching
upper bound.


\begin{cor}[\textbf{EXPSPACE-completeness of structural liveness}]\label{cor:conservEScomplete}
\hfill
	\begin{enumerate}
		\item			
	Structural liveness for conservative nets is EXPSPACE-complete.
\item
	Structural liveness for structurally bounded nets is EXPSPACE-complete.
	\end{enumerate}
	\end{cor}	
\begin{proof}
	Since EXPSPACE$=$NEXPSPACE,
it suffices to provide a nondeterministic algorithm working in
	exponential space
	that verifies positive instances of the considered problems.
	By Lemma~\ref{lem:consstrbound}, the positive instances of
	both problems coincide; in particular, every structurally live
	structurally bounded net is conservative. 

	Given a structurally live conservative $d$-dim net $A$, an 
	algorithm can guess a live
	marking $x\in\setN^d$ as guaranteed by Theorem~\ref{th:upper}.
 Since  $\nnorm{x}\leq f(\nnorm{A},d)$ for the corresponding 2-exp function
	$f$
	satisfying	$f(m,i)\leq(2+m)^{2^{\text{poly}(i)}}$, the
	binary representation of $x$ has size at most
	$O(2^{\text{poly}(d)}\cdot\log (2+\nnorm{A}))$, which fits in exponential space.
Finally, verifying the liveness of $x$ for $A$ can be performed in
	polynomial space with respect to
	the binary encoding of $A$ and $x$ (see, \eg,~\cite[Lemma
	5]{DBLP:conf/apn/MayrW14}), which remains within exponential
	space relative to the original input size. 
\end{proof}
